\documentclass[11pt]{article}

\usepackage[a4paper,margin=1in]{geometry}
\usepackage{amsmath,amssymb}
\usepackage[hidelinks]{hyperref}
\usepackage{enumitem}

\emergencystretch=4em
\sloppy

\title{Follow-up Review on Finding F3 for\\
\texttt{rational\_cut\_and\_project\_gap\_periods.tex}}
\author{Codex Review Notes}
\date{May 13, 2026}

\begin{document}

\maketitle

\section*{Executive Summary}

I reviewed the partial implementation of Finding~F3 in the irrational
aperiodicity proof, focusing on the passage where a finite projected gap
period is lifted to a lattice translation of the accepted set.

The main mathematical concern behind F3 has been addressed: the text no
longer speaks of \(-\sigma\) as a ``period''.  It now correctly describes
\(-\sigma\) as a translation of the projected set, derived from
\(\Lambda+\sigma=\Lambda\).  This removes the potentially misleading
formulation from the earlier version.

The remaining question is mostly stylistic and about alignment with the Lean
proof: the paper still proves the stronger statement \(A+v=A\), whereas the
Lean formalisation only needs and proves the forward inclusion
\[
  A+v\subseteq A.
\]
This is not a correctness blocker, but the proof can still be simplified if
desired.

\section*{Current State}

The relevant paper passage is now around
\[
\texttt{LaTeX/rational\_cut\_and\_project\_gap\_periods.tex:1359--1375}.
\]
It argues:
\[
  \ell_{i+\lambda}=\ell_i+\sigma
  \quad\Longrightarrow\quad
  \Lambda+\sigma=\Lambda,
\]
then chooses \(v\in\mathbb Z^2\) with \(\tilde p(v)=\sigma\), proves
\[
  A+v\subseteq A,
\]
and then derives the reverse direction by using
\[
  \Lambda-\sigma=\Lambda.
\]
The text now says explicitly that \(\Lambda+\sigma=\Lambda\) is equivalent
to \(\Lambda-\sigma=\Lambda\).  This is mathematically fine.

The Lean proof, by contrast, has the sharper shape:
\[
\texttt{period\_lifts\_to\_lattice\_translation}
\quad\Rightarrow\quad
\forall z\in A,\ z+v\in A,
\]
see \texttt{Irrational.lean:704--709}.  Then
\texttt{lattice\_translation\_must\_be\_zero} uses only that forward
inclusion, see \texttt{Irrational.lean:798--801}.

\section*{Remaining Points}

\begin{enumerate}[label=\textbf{R\arabic*.},leftmargin=*]

\item \textbf{Optional: align the paper proof with the Lean proof.}

The paper still states
\[
  A+v=A
\]
and Step~3 uses this equality as ``invariance''.  This is stronger than
necessary.  The Lean proof shows that forward invariance alone gives the
contradiction: if \(s(v)=\tau>0\), choose an internal point \(t\) near the
upper endpoint of \(W\), so \(t+\tau\notin W\); if \(\tau<0\), choose \(t\)
near the lower endpoint.

\textbf{Importance: low to medium.}  This is not a mathematical flaw.  The
current paper proof is valid as written, assuming the short reverse-translation
argument is accepted.  Aligning with Lean would make the paper shorter, closer
to the formalisation, and less likely to invite questions about the reverse
translation step.

\item \textbf{Optional: make the reverse-translation argument fully explicit
if keeping \(A+v=A\).}

The current wording says that applying the same argument with translation
\(-\sigma\) gives \(z-v\in A\).  This is now much better than the old
``period \(-\sigma\)'' wording.  For maximum referee-proof clarity, one could
spell out the one-line calculation:
\[
  \tilde p(z-v)=\tilde p(z)-\sigma.
\]
Since \(\Lambda-\sigma=\Lambda\), the value \(\tilde p(z)-\sigma\) is
attained by some accepted point; injectivity of \(\tilde p\) then forces that
point to be \(z-v\).

\textbf{Importance: low.}  The present argument is understandable.  This
would only remove one small ``same argument'' compression.

\item \textbf{Optional: update Step~3 wording if the proof is simplified.}

If the paper is changed to use only \(A+v\subseteq A\), then Step~3 should no
longer say that \(A+v=A\) translates to
\[
  (W\cap s(\mathbb Z^2))+\tau
  =
  W\cap s(\mathbb Z^2).
\]
It should instead use the forward implication
\[
  t\in W\cap s(\mathbb Z^2)
  \quad\Longrightarrow\quad
  t+\tau\in W.
\]
This is exactly the form used by Lean.

\textbf{Importance: conditional.}  If the paper keeps the current
\(A+v=A\) proof, no change is needed here.  If the paper is aligned with Lean,
this wording must be adjusted.

\end{enumerate}

\section*{Assessment}

Finding~F3 should now be considered \emph{mostly resolved}.  The original
problematic phrase has been removed, and the current proof is mathematically
sound.  The only remaining work is optional polish:
\begin{itemize}[leftmargin=*]
  \item either keep \(A+v=A\) and expand the reverse-translation argument by
        one explicit sentence;
  \item or simplify the paper to match Lean by proving only
        \(A+v\subseteq A\) and using the endpoint contradiction directly.
\end{itemize}

For arXiv submission, I would not treat the remaining F3 items as blockers.
If there is time, I slightly prefer the Lean-aligned forward-inclusion version,
because it is shorter and conceptually cleaner.  But the current version is
acceptable from a mathematical correctness standpoint.

\section*{Checks Performed}

\begin{itemize}[leftmargin=*]
  \item I inspected the updated irrational proof around
        lines 1359--1395 of the main paper.
  \item I compared it with the Lean declarations
        \texttt{period\_lifts\_to\_lattice\_translation} and
        \texttt{lattice\_translation\_must\_be\_zero} in
        \texttt{Irrational.lean}.
  \item I searched the paper for the old problematic wording; no occurrence
        of ``period \(-\sigma\)'' remains.
\end{itemize}

\end{document}
